\chapter{Назначение и текущая область применения}

\program{} --- Python-пакет для построения воспроизводимой фотометрической базы по архивным астрономическим наблюдениям. Основной принцип проекта --- хранить не только конечную кривую блеска, но и provenance: какие кадры, каталог, калибровочные поколения, quality gates и программная версия породили результат.

\section{Три уровня рабочего процесса}

Проект удобно рассматривать как три последовательные стадии:

\begin{enumerate}
  \item \textbf{Подготовка изображения.} Получение raw FITS, выбор bias/dark/flat, редукция, нормализация заголовков и получение рабочего FITS. Полноценный production-конвейер этой стадии ещё не реализован.
  \item \textbf{Фотометрия одной серии.} Работа с набором уже подготовленных FITS одной серии: каталог, детекция, forced photometry, aperture/ZP/ensemble/color calibration, QC и science measurements. Это наиболее реализованная, протестированная и оптимизированная часть программы.
  \item \textbf{История объекта по многим сериям.} Поиск всех наблюдений одного объекта, межсерийная калибровка, построение долговременных кривых блеска и последующий variability analysis. Это следующий основной этап разработки.
\end{enumerate}

\section{Поддерживаемые данные}

Текущая кодовая база содержит instrument profiles и рабочие ветви для MASTER и RoboPhot. Конкретная доступность каталожной, reference/color и detector-specific логики зависит от профиля инструмента и фильтра.

\section{Что не следует ожидать от текущей версии}

\begin{itemize}
  \item Автоматическая bias/dark/flat редукция raw архива пока не является завершённой пользовательской стадией.
  \item Единый production-продукт «объект по всем сериям» ещё не завершён.
  \item Некоторые residual-noise quantities сохраняются как диагностика и не должны автоматически трактоваться как применённая поправка.
\end{itemize}
